\documentclass[12pt]{jsarticle}
\usepackage{amsmath,amssymb}
\pagestyle{empty}
\begin{document}
\setlength{\baselineskip}{22pt}

{\huge 大阪大学　理系数学　2013}

\newpage

\begin{flushleft}
{\huge 1}
　三角関数の極限に関する公式
\[ \lim_{x\to0}\frac{\sin x}{x}=1 \]
を示すことにより，$\sin x$の導関数が$\cos x$であることを証明せよ．
\end{flushleft}

\newpage

\begin{flushleft}
{\huge 2}
　不等式
\[ 1\leqq||x|-2|+||y|-2|\leqq3 \]
の表す領域を$xy$平面上に図示せよ．
\end{flushleft}

\newpage

\begin{flushleft}
{\huge 3}
　4個の整数
\[ n+1, \hspace{1zw} n^3+3, \hspace{1zw} n^5+5, \hspace{1zw} n^7+7 \]
がすべて素数となるような正の整数$n$は存在しない．これを証明せよ．
\end{flushleft}

\newpage

\begin{flushleft}
{\huge 4}
　$xyz$空間内の3点$O(0, \, 0, \, 0)$，$A(1, \, 0, \, 0)$，$B(1, \, 1, \, 0)$を頂点とする三角形$OAB$を
$x$軸のまわりに1回転させてできる円すいを$V$とする．
円すい$V$を$y$軸のまわりに1回転させてできる立体の体積を求めよ．
\end{flushleft}

\newpage

\begin{flushleft}
{\huge 5}
　$n$を3以上の整数とする．$n$個の球$K_1, \, K_2, \, \cdots\cdots, \, K_n$と
$n$個の空（から）の箱$H_1, \, H_2, \, \cdots\cdots, \, H_n$がある．
以下のように，$K_1, \, K_2, \, \cdots\cdots, \, K_n$の順番に，球を箱に1つずつ入れていく．\\
　まず，球$K_1$を箱$H_1, \, H_2, \, \cdots\cdots, \, H_n$のどれか1つに無作為に入れる．
次に，球$K_2$を，箱$H_2$が空ならば箱$H_2$に入れ，
箱$H_2$が空でなければ残りの$n-1$個の空の箱のどれか1つに無作為に入れる．\\
　一般に，$i=2, \, 3, \, \cdots\cdots, \, n$について，
球$K_i$を，箱$H_i$が空ならば箱$H_i$に入れ，
箱$H_i$が空でなければ残りの$n-i+1$個の空の箱のどれか1つに無作為に入れる．
\begin{description}
\item[(1)]$K_n$が入（はい）る箱は$H_1$または$H_n$である．これを証明せよ．
\item[(2)]$K_{n-1}$が$H_{n-1}$に入る確率を求めよ．
\end{description}
\end{flushleft}

\end{document}
